%! TeX root = ../phd-thesis.tex
%! Author = davidedomini

%----------------------------------------------------------------------------------------
\chapter{Conclusions and Future Work}
\label{chap:wrapup}
\minitoc
%----------------------------------------------------------------------------------------

This thesis investigated how cooperative learning can be engineered
 in \glspl{CAS} under conditions that complicate conventional learning methods.
%
The starting point was that many-agent deployments invalidate several assumptions 
 commonly made by centralized learning: 
 data cannot always be pooled, global state may be unavailable, 
 the population and communication topology may change, 
 and individual devices may lack the resources required to train or execute complex models locally.
%
Addressing these conditions ultimately requires attention to how learning,
 collective organization, and deployment interact.
%
This thesis examined parts of that broader problem in separate experimental settings,
 developing methods and evidence for each without combining them into a single system.

The contributions followed three complementary directions.
%
\Cref{chap:selforg-fl} studied how federated-learning groups can emerge 
 and reorganize through local field-based coordination.
%
\Cref{chap:scalable-rl} examined how reinforcement learning can exploit neighborhood exchanges
 and graph-based policies to remain decentralized and operate across different population sizes.
%
\Cref{chap:tasksplacement} examined the placement of distributed applications,
 combining heterogeneous graph learning with aggregate information to support component-level offloading 
 across the edge--cloud continuum.
%
This final chapter consolidates the evidence obtained from these contributions,
 answers the research questions introduced in~\Cref{chap:introduction}, 
 and identifies the main directions required to move from controlled evaluations 
 to long-lived operational systems.


%=============================================================================
\section{Conclusions}
\label{sec:conclusions}
%=============================================================================

Across the investigated settings, the central finding is that information about
 collective structure can improve learning and decentralized decisions
 without requiring each device to observe or control the entire system.
%
The relevant structure differed across the contributions:
 learning groups in federated learning, interaction neighborhoods in swarm coordination,
 and shared infrastructure conditions in application offloading.
%
Field-based coordination organized collaboration and computed collective information,
 while neighborhood exchanges and graph-based policies incorporated relationships
 among agents into learning and execution.
%

\paragraph{RQ1: How can privacy-aware cooperative learning be organized among
many distributed devices without centralizing raw data or relying on a fixed
central coordinator?}
%
The results of~\Cref{chap:selforg-fl} show that field-based coordination 
 can construct the organizational structure required by \gls{FL}
 directly from local interactions.
%
FBFL replaced a single global federation 
 with a set of self-organizing spatial regions, 
 each temporarily coordinated by an elected leader. 
% 
This organization realized a form of \gls{CFL} 
 in which each region trained and maintained its own model, 
 allowing multiple models specialized to different local data distributions to coexist. 
% 
When leaders failed, 
 the regions were reconstructed and training continued 
 without external reconfiguration or lasting accuracy loss.
%
PSFL generalized this organization by incorporating model compatibility 
 into the distance used for region formation. 
% 
This allowed the number and membership of federations 
 to emerge during execution instead of being fixed in advance, 
 and enabled devices to revise their membership as they moved.

These results establish that the control plane of \gls{FL} need not be centralized 
 or statically configured.
%
Keeping training data on the originating devices 
 and restricting model exchange to dynamically formed neighborhoods reduces the need 
 for global data collection and limits the scope of collaboration.
%
This should nevertheless be interpreted as a privacy-aware architecture 
 rather than a formal privacy guarantee. 
% 
Protecting the information contained in model parameters 
 and updates constitutes a complementary line of research in \gls{FL}, 
 largely orthogonal to the organizational problem addressed here. 
% 
Mechanisms such as secure aggregation 
 and differential privacy can be integrated at the model-exchange 
 and aggregation layers without altering the self-organizing principles 
 used to form and maintain federations. 
% 
The proposed methods therefore provide an organizational substrate 
 on top of which stronger privacy guarantees can be implemented, 
 rather than replacing the mechanisms that provide those guarantees.

\paragraph{RQ2: How can cooperative learning account for spatially heterogeneous
data and adapt to temporal distribution shifts caused by mobility and changing
environments?}
FBFL demonstrated that spatial structure can be used as an inductive bias for
collaboration when proximity correlates with the data-generating process. Under
spatially partitioned non-\gls{IID} data, regional models avoided interference
between incompatible distributions and substantially outperformed methods that
forced all devices to optimize a single global model.
%
However, 
 FBFL relied exclusively on physical proximity to cluster devices. 
% 
This assumption is effective when spatial closeness reflects similarity 
 between local data distributions, 
 but the strength and spatial scale of this relation may be difficult 
 to determine a priori and may vary across the environment. 
% 
PSFL relaxed this assumption by replacing the purely spatial distance 
 with a configurable metric incorporating loss-based model compatibility. 
% 
Regions could therefore expand where devices learned compatible models 
 and contract or separate where their models diverged, 
 allowing the emerging federations to follow the statistical structure of the data more closely 
 while retaining local communication.

Mobility revealed a distinction between organizational and model adaptation.
%
The field-based structure can track changes in device position 
 and reconstruct appropriate federations, 
 but changing collaborators does not preserve knowledge acquired in regions visited previously.
%
For a mobile device, 
 spatial heterogeneity becomes a temporal sequence of distributions 
 and can therefore produce catastrophic forgetting.
%
The preliminary C2FL results show that local replay and dwell-time-aware model integration 
 can mitigate this effect while retaining the benefit of the current regional federation.


\paragraph{RQ3: How can decentralized learning remain effective and scalable
when devices have access only to local observations, communicate with a limited
neighborhood, and operate in collectives whose size may change?}
\Cref{chap:scalable-rl} investigated two distinct dimensions of decentralized scalability:
 where learning information is exchanged and whether the learned policy 
 can be executed by populations of different sizes.
%
The neighbor-based training experiments showed that full independence 
 is not the only alternative to centralized training.
%
Exchanging experiences or averaging model information within bounded neighborhoods 
improved the stability and performance of decentralized deep Q-learning, 
approaching centralized training in the evaluated coordination task.
%
The comparison also exposed the cost of different exchanges. 
%
Experience sharing provided a favorable balance because its communication footprint remained small,
 whereas exchanging complete model information incurred a larger cost that grew with neighborhood size.
% 
Local cooperation can therefore recover part of the benefit of centralized training, 
 but the exchanged representation determines whether that improvement remains feasible.

DGQL addressed population scaling at the policy level by representing agents 
 and their interactions as a graph and sharing the same graph-based policy across the collective.
%
The learned policy remained functional when deployed without retraining 
 in swarms larger than the training population: 
 agents continued to approach their goals and maintained obstacle avoidance.
%
Performance, however, 
 degraded as population size diverged from the training configuration. 
% 
Increased density changed the local interaction patterns observed by the policy 
 and exposed the limited receptive field of a fixed neighborhood 
 and shallow message-passing architecture. 
% 
In the obstacle-avoidance task, this shift appeared as a safety--efficiency trade-off:
 collisions remained rare, but agents became increasingly conservative and less effective at reaching the goal.

\paragraph{RQ4: How can learning-based decisions be deployed across heterogeneous
and resource-constrained edge--cloud infrastructures while accounting for both
individual device constraints and collective system-level phenomena?}
IDHGQL addressed this problem by combining three representations.
%
The macro-component model decomposed applications into selectively placeable
 units while preserving their data and execution dependencies.
%
The heterogeneous graph distinguished application devices,
 edge and cloud resources, communication links, and workflow relations.
%
Aggregate computing enriched each local observation with self-stabilizing
 summaries of collective conditions that could not be inferred reliably from the
 state of one device alone.

The evaluation showed that the resulting policy could express different
 trade-offs among battery preservation, infrastructure cost, and server load.
%
More importantly, the density-aware experiments isolated the contribution of
 collective information. 
% 
With aggregate density included in the observation,
 devices learned to execute components locally in congested areas and offload from
 sparser regions; when this information was removed, the agents converged toward a
 largely uniform strategy that did not reflect their position in the collective.
%
The learned policy also adjusted its decisions when devices moved between dense
 and sparse areas.

\paragraph{Cross-cutting findings.}
Taken together, the four answers identify three principles that recur across the
contributions.
%
First, the collaboration topology is part of the learning problem. Choosing who
exchanges models, experiences, messages, or infrastructure information changes
both the statistical objective and the cost of reaching it. This topology should
therefore be adaptive and task-aware rather than treated as fixed communication
plumbing.
%
Second, locality is not only a limitation. Spatial and graph neighborhoods reduce
the information available to an individual device, but they also provide structure
that can constrain collaboration, bound communication, and expose the context in
which local decisions take effect.
%
Third, scalability is multidimensional. Removing a central server, keeping model
size independent of population size, or supporting variable graphs addresses only
one source of growth. Communication volume, behavioral distribution shift,
resource contention, and adaptation time must be considered jointly.

In the federated learning and offloading studies, macroprogramming provided
 a computational layer for constructing learning groups and collective state
 through local interactions. The reinforcement learning studies addressed
 complementary questions through neighborhood exchanges and graph-based policies.




%=============================================================================
\section{Future Work}
\label{sec:future-work}
%=============================================================================

The thesis establishes mechanisms for self-organizing learning,
 scalable policy representation, and collective-aware deployment, 
 each investigated through focused experimental settings designed to isolate 
 the corresponding engineering challenges.
%
Building on these results, 
 future research should strengthen adaptation within each layer and examine how the proposed mechanisms 
 interact when integrated into long-running systems.

\subsection{Autonomous and Continual Organization}

The formation of learning groups should become less dependent 
 on the task model and on manually selected thresholds.
%
PSFL estimates compatibility by exchanging 
 and evaluating the models used for the primary task. 
% 
Although effective, 
 this couples federation discovery to model architecture 
 and incurs repeated communication and inference costs.
%
Lightweight auxiliary representations, 
 such as novelty estimators based on random network distillation, 
 offer a way to identify distributional compatibility before or independently of task training. 
% 
Future work should determine when this discovery process must be repeated 
 and replace periodic clustering with event-driven mechanisms 
 triggered by statistically meaningful changes.
%
Thresholds governing compatibility, federation radius, 
 and leader selection should likewise adapt to local uncertainty 
 and observed stability rather than remain fixed for an entire deployment.

Long-lived mobility also requires continual adaptation of model knowledge.
%
The initial C2FL results motivate a broader comparison among replay, regularization, 
parameter isolation, and hybrid continual-learning strategies in decentralized federations. 
%
Replay buffers must be bounded on resource-constrained devices, 
 making sample selection and replacement policies part of the system design. 
% 
Adaptation should be evaluated under gradual shifts, recurring regions, irregular mobility, 
 and delayed federation updates rather than only under abrupt transitions. 
% 
A further challenge is to distinguish genuine distribution change from temporary noise 
 so that federation reorganization and model consolidation are not triggered unnecessarily.

\subsection{Intrinsic Model Similarity}

Complementary to auxiliary representations, 
 compatibility could be estimated directly from intrinsic properties of the trained task models, 
 without evaluating them on local data. 
% 
This requires more than applying a conventional vector similarity measure, 
 such as cosine similarity, to flattened parameters. 
% 
Neural networks admit multiple parameterizations of the same function: 
 for example, hidden units can be permuted together with their incident connections 
 without changing the network output, despite producing substantially different parameter vectors. 
% 
A useful model-similarity function must therefore be invariant, 
 or at least robust, to such function-preserving transformations.
%
Developing this function would also provide a basis for studying how model trajectories diverge 
 when training data are drawn from homogeneous or heterogeneous distributions, 
 and for identifying which structural properties of the models reliably indicate statistical compatibility.

\subsection{Scale- and Resource-Aware Learning}

Both the federated and reinforcement learning studies show that 
 fixed local interaction parameters become limiting when density, topology, 
 or resource availability changes.
%
Neighborhood size, communication frequency, and the type of exchanged information 
 should be selected at runtime according to their expected benefit.
%
Agents could communicate when uncertainty is high, 
 when a neighbor holds novel experience, or when local estimates indicate that 
 the collective is becoming inconsistent. 
% 
Such event-driven protocols would make communication a controlled
 resource rather than a fixed side effect of every training round.

Graph-based policies similarly require explicit awareness of scale.
%
Training across randomized population sizes and densities, conditioning policies on local density, 
 and progressively increasing coordination complexity may reduce the distribution shift observed by DGQL. 
% 
Adaptive graph construction could vary neighborhood size and receptive field with the task, 
 while hierarchical or attention-based message passing could combine local responsiveness with
 longer-range information. 
% 
These extensions must be evaluated against their communication and inference costs: 
 increasing message-passing depth may improve coordination while undermining the bounded local interaction 
 that makes the architecture deployable.

Resource adaptation should also extend beyond a uniform model shared by all devices. 
%
SParSeFuL showed that sparsification affects not only computational cost and accuracy, 
 but also the similarity estimates from which federations emerge.
%
Future methods should therefore adapt model capacity jointly with federation membership, 
 accounting for heterogeneous processors, memory budgets, battery levels, and link quality. 
% 
This requires metrics that include transmitted bytes, training energy, 
 and wall-clock adaptation time in addition to predictive performance.

\subsection{Joint Learning and Deployment}

The current contributions optimize learning organization 
 and application placement in separate stages. 
% 
Operational systems will require these decisions to interact.
%
The device hosting a model, 
 the neighbors with which it exchanges information,
 and the infrastructure on which its components execute all influence latency,
 energy consumption, and the freshness of the available knowledge. 
% 
A natural extension is therefore a unified control layer that 
 can jointly adapt federation membership, model capacity, communication, and component placement.

Such integration should preserve a separation between different timescales.
%
Fast field computations can track topology, congestion, and failures; 
 learning can optimize decisions whose consequences unfold over longer horizons; 
 and continual-learning mechanisms can consolidate knowledge across persistent changes. 
% 
Allowing one mechanism to control all three timescales directly 
 would make the system difficult to stabilize and diagnose. 
% 
Explicit interfaces between coordination fields, learned policies, 
 and deployment actions would instead make their dependencies observable 
 and allow each layer to adapt at an appropriate rate.

For IDHGQL, 
 the immediate next step is to coordinate the placement of dependent
 components rather than letting each agent choose independently. 
% 
Agents may need to negotiate compatible placements, reserve shared capacity, 
 or account for the future decisions of nearby devices. 
% 
Multi-objective policies should also move beyond fixed linear scalarization, 
 either by learning a set of Pareto-efficient policies or by adapting priorities to the operating context. 
% 
These extensions would turn collective-aware individual offloading into coordinated deployment 
 of distributed workflows.

\subsection{Validation in Operational Systems}

The experiments in this thesis deliberately isolate individual mechanisms in simulation. 
%
Establishing external validity requires deployments in which sensing, communication, computation, 
 and mobility introduce delays and failures that are only approximated by the current models.
%
Future evaluations should include physical robot or sensor collectives 
 and edge--cloud testbeds with heterogeneous hardware, intermittent connectivity,
 asynchronous execution, packet loss, and devices joining or leaving during operation. 
% 
Population size and topology should vary during both training 
 and deployment, rather than being changed only between experimental runs.

Evaluation criteria must also broaden beyond final task accuracy or reward.
%
Relevant measures include recovery time after disruption, communication volume,
 energy consumed by training and inference, stability during reorganization,
 fairness across devices and regions, and the quality of service experienced 
 while adaptation is still in progress. 
% 
Privacy claims should be assessed against explicit adversary models, 
 including information leakage from exchanged parameters and collective summaries. 
% 
Finally, 
 comparisons should include the engineering cost of configuration: 
 an approach that performs well only after selecting the correct number of clusters, 
 neighborhood size, reward weights, or aggregation interval may be less useful 
 in an open system than one with slightly lower peak performance but reliable self-configuration.

Progress along these directions would move the proposed methods from separate proofs of feasibility 
 toward a unified engineering discipline for cooperative learning in \glspl{CAS}. 
% 
The defining challenge is not merely to make each device learn, 
 but to ensure that learning, coordination, and deployment continue 
 to support one another as the collective changes. 
%
The evidence developed in this thesis indicates that this objective is attainable 
 when collective structure is made explicit and treated as a first-class element of the learning system.